\section{Background}
\label{ch:background}

This section develops the background the thesis relies on: the energy-based view of sequence generation, Langevin dynamics, the Metropolis--Hastings correction, the two concrete samplers studied in the experiments, score matching and the concrete score, and GFlowNets as an amortized alternative. The aim throughout is to expose the assumption each tool makes rather than to survey the tools exhaustively.

\subsection{Autoregressive Language Models as Energy Functions}
\label{sec:bg-ebm}

An autoregressive language model defines a probability distribution over a sequence of tokens $\bm{x} = (x_1, \dots, x_T)$ by the chain rule of probability,
\begin{equation}
    \log p_{\text{LM}}(\bm{x}) = \sum_{t=1}^{T} \log p_{\text{LM}}(x_t \mid x_{<t}),
    \label{eq:bg-chain}
\end{equation}
where each conditional $p_{\text{LM}}(x_t \mid x_{<t})$ is produced by a neural network, here a Transformer \citep{vaswani2017attention}, applied to the prefix: it reads the prefix, produces logits over the vocabulary, and a softmax turns them into the next-token distribution. The model is trained by \textbf{teacher forcing}, maximizing \eqref{eq:bg-chain} on a large corpus with the true prefix supplied at every position, which decomposes into a sum of independent next-token prediction problems.

It is worth being precise about what this training objective does and does not shape, because Hypothesis B of Section~\ref{sec:intro-ebm} is stated in exactly these terms. The objective shapes each conditional $p_{\text{LM}}(x_t \mid x_{<t})$ to be an accurate predictor of the next token given a \emph{correct} history drawn from the data distribution, and modern models optimize it extraordinarily well. What it never does is ask the joint quantity $\log p_{\text{LM}}(\bm{x})$ to behave in any particular way as $\bm{x}$ is varied away from the data, for instance by substituting one token for another in the middle of an otherwise fixed sequence. The model is never trained on such counterfactual sequences, never asked to assign them calibrated scores, and never penalized for the shape of the surface it induces over them, so nothing internal to the objective makes the joint log-likelihood smooth, or well-calibrated as a measure of quality, or possessed of a gradient pointing toward better sequences. That is a real gap, and it makes Hypothesis B the natural first explanation of a failure; Section~\ref{ch:results} nonetheless finds it too strong, since the quantity a local proposal needs turns out to be present in the trained conditionals and reachable without any change of objective.

With that caveat noted rather than resolved, the energy-based reformulation is immediate. Identifying the Boltzmann form of \eqref{eq:intro-ebm} with the language model, one simply sets
\begin{equation}
    \energy_{\text{LM}}(\bm{x}) = -\log p_{\text{LM}}(\bm{x}),
    \label{eq:bg-energy}
\end{equation}
so that sequences of high probability are sequences of low energy. Sampling text from the language model is then, formally, sampling from the Boltzmann distribution of this energy. For controllable generation an additive constraint $\lambda\, C(\bm{x})$ is appended as in \eqref{eq:intro-energy}, where $C$ is typically the negative log-probability an attribute classifier assigns to the desired label, so low energy means both fluent and on-attribute. The central computational question, to which the next two sections are devoted, is how one draws such samples at all.

\subsection{Langevin Dynamics}
\label{sec:bg-langevin}

Suppose, for a moment, that the state $\bm{s}$ lived in a continuous space $\R^D$ and that the energy were differentiable everywhere. Then one could sample from the distribution $p(\bm{s}) \propto \exp(-\energy(\bm{s}))$ using \textbf{Langevin dynamics}, a gradient-guided Markov chain Monte Carlo method with a long history in statistics and physics \citep{roberts1996exponential, welling2011bayesian}. The discretized update, which is what one actually runs, is
\begin{equation}
    \bm{s}_{t+1} = \bm{s}_t + \frac{\epsilon_t}{2}\,\grad_{\bm{s}} \log p(\bm{s}_t) + \sqrt{\epsilon_t}\,\bm{\xi}_t, \qquad \bm{\xi}_t \sim \mathcal{N}(\bm{0}, \bm{I}),
    \label{eq:bg-langevin}
\end{equation}
where $\epsilon_t$ is a step size at iteration $t$. The update has two parts, and their interplay is what distinguishes sampling from optimization.

The first part, the \textbf{drift}, is a step of size $\epsilon_t/2$ along the gradient of the log-density. Iterated alone it is gradient ascent on $\log p$ and drives the state monotonically toward the nearest mode, which would suffice if the goal were to \emph{find} the most probable point; but the goal is a \emph{representative sample}, and a distribution is more than its mode. The \textbf{diffusion} term injects fresh Gaussian noise at every step, scaled by $\sqrt{\epsilon_t}$, which prevents the chain from collapsing onto the mode: it knocks the state off the deterministic ascent path, letting it climb away from a mode, cross into a neighbouring region, and in the limit visit every part of the distribution with the correct frequency. The drift acts like gravity pulling a marble downhill on a landscape whose valleys are the high-probability regions, while the diffusion vibrates the surface so the marble does not settle in the first valley it finds. \textbf{Annealing} the step size toward zero on a suitable schedule makes the distribution of $\bm{s}_t$ converge to the target $p(\bm{s})$ rather than to its mode \citep{welling2011bayesian}, but that fair-draw reading holds only under the technical conditions of the result, in particular a suitable schedule and a well-behaved landscape, and does not follow in general simply because the vibration was switched off.

The guarantees behind the method rest on regularity conditions that are easy to pass over and matter a great deal later. The discretization error of \eqref{eq:bg-langevin} grows with the step size, so the finite-step chain is an approximate sampler whose bias increases with $\epsilon_t$ \citep{roberts1996exponential}; the Metropolis-adjusted variant of the next section removes that bias in principle, but its acceptance rate degrades sharply once the drift loses smoothness, because the correction depends on the reverse move being predictable from the forward one. Non-asymptotic guarantees are established under smoothness of the target and typically under log-concavity or a related condition \citep{dalalyan2017theoretical, durmus2017nonasymptotic}, and in particular require the drift to be \textbf{Lipschitz continuous} \citep{roberts1996exponential}, so that a small step lands where the gradient at the starting point predicted. The continuous sampler of Section~\ref{sec:bg-samplers} violates these assumptions, because the drift it uses is computed from a differentiable pathway whose prediction target changes discontinuously at each Voronoi cell boundary. Section~\ref{sec:results-mh} measures the consequence.

\subsection{The Metropolis--Hastings Correction}
\label{sec:bg-mh}

Any finite step size in \eqref{eq:bg-langevin} introduces a bias: the chain samples not from $p(\bm{s})$ but from a perturbed distribution that depends on the step size. That bias can be removed exactly by the \textbf{Metropolis--Hastings} (MH) correction \citep{metropolis1953equation, hastings1970monte}, which turns a biased proposal into an exact sampler by accepting or rejecting each move. Having proposed a move from $\bm{s}$ to $\bm{s}'$ according to a density $q(\bm{s}' \mid \bm{s})$, one accepts with probability
\begin{equation}
    \alpha(\bm{s}' \mid \bm{s}) = \min\left\{ 1,\; \frac{p(\bm{s}')\, q(\bm{s} \mid \bm{s}')}{p(\bm{s})\, q(\bm{s}' \mid \bm{s})} \right\},
    \label{eq:bg-mh}
\end{equation}
and otherwise rejects it and remains at $\bm{s}$, counting the current state again. When Langevin dynamics is equipped with this correction the method is known as the Metropolis-adjusted Langevin algorithm, or MALA.

The acceptance ratio factors into two parts. The \textbf{target ratio} $p(\bm{s}')/p(\bm{s})$ compares the proposed state to the current one under the distribution being sampled and expresses the sampler's preference for good states. The \textbf{proposal ratio} $q(\bm{s} \mid \bm{s}')/q(\bm{s}' \mid \bm{s})$ compares how likely the mechanism was to suggest the reverse move against the forward one, and it is what enforces \textbf{detailed balance}, penalizing moves that are easy to make but hard to undo, which if accepted freely would let the chain drift and distort the distribution.

Two facts about the correction recur in the results. First, a great many energy-based decoders omit it entirely, and without it the noise-augmented gradient update of \eqref{eq:bg-langevin} is not a sampler but a biased optimizer with a particular kind of stochasticity, so calling its output a sample from the energy is strictly incorrect. The distinction is not pedantic, because a biased optimizer stopped early and a correct sampler run to convergence can produce very different outputs; how far it reframes any published system is a separate question, and one this thesis can answer only for the implementations it runs itself. Second, the proposal ratio is delicate to compute for a Langevin proposal, because the reverse density $q(\bm{s} \mid \bm{s}')$ must be evaluated from the drift \emph{at the proposed point}. If that drift differs wildly from the drift at $\bm{s}$, which is what happens when the two points fall on opposite sides of a discontinuity, the reverse proposal places the original state deep in the tail of a Gaussian centred far away, driving $q(\bm{s} \mid \bm{s}')$ and the whole acceptance probability toward zero. Section~\ref{sec:results-mh} decomposes the ratio and reports which term is responsible.

\subsection{Sampling in a Discrete Space: DLS and CLS}
\label{sec:bg-samplers}

The difficulty that motivates the entire thesis is that text is not continuous. A sequence is a list of discrete tokens drawn from a finite vocabulary $V$, and there is no state \emph{between} two tokens, so Langevin dynamics, which assumes a continuous differentiable space in which infinitesimal steps are meaningful, cannot be applied without modification. Two strategies for adapting it are studied here, and they differ precisely in where they confront the discreteness, which is why studying both is what allows the results to be attributed to the energy rather than to one sampler's idiosyncrasies.

The \textbf{Discrete Langevin Sampler} (DLS), following the discrete gradient-based samplers of \citet{grathwohl2021oops} and \citet{zhang2022langevin}, never leaves token space: its state is always a genuine sequence of tokens, and it uses the gradient of the log-likelihood with respect to the continuous input embedding of a masked position only as a \emph{score} for building a categorical proposal over which token to place there next. Let $\ve(x_i)$ be the embedding of the current token at masked position $i$ and $\vg = \grad_{\ve_i}\log p_{\text{LM}}(\bm{x})$ the gradient with respect to it. For a candidate token $v$ the proposal combines a term penalizing the displacement $\norm{\ve(v) - \ve(x_i)}$, which keeps proposals local, with a term proportional to $\tfrac{1}{2}\vg^\top(\ve(v) - \ve(x_i))$, which biases the proposal toward tokens the gradient suggests will lower the energy. That second term is a \textbf{first-order Taylor surrogate} for the true change in log-likelihood a substitution would produce.

Written out in full, the discrete sampler adapts the continuous Langevin proposal of \eqref{eq:bg-langevin} to a categorical distribution over the vocabulary by the construction of \citet{zhang2022langevin}. A Gaussian proposal centred at $\bm{s} + \tfrac{\alpha}{2}\grad\log p(\bm{s})$ assigns a candidate $\bm{s}'$ a log-density proportional to $-\tfrac{1}{2\alpha}\norm{\bm{s}' - \bm{s} - \tfrac{\alpha}{2}\grad\log p(\bm{s})}^2$, which on expanding and discarding terms constant in $\bm{s}'$ leaves a term linear in the displacement plus a quadratic penalty on the displacement itself. Transplanting this to the discrete setting, where the candidate states are token embeddings $\ve(v)$ and the current state is $\ve(x_i)$, gives a categorical proposal whose logit for token $v$ is
\begin{equation}
    \log q(v \mid x_i) \propto \tfrac{1}{2}\,\vg^\top\bigl(\ve(v) - \ve(x_i)\bigr) - \tfrac{1}{2\alpha}\norm{\ve(v) - \ve(x_i)}^2,
    \label{eq:bg-discrete-proposal}
\end{equation}
where $\vg = \grad_{\ve_i}\log p_{\text{LM}}(\bm{x})$. The first term, the \textbf{gradient-alignment term}, is the surrogate: it rewards candidates whose displacement from the current embedding aligns with the gradient, and it is precisely the first-order Taylor estimate of how much the log-likelihood would change if the substitution were made, since to first order $\log p(\bm{x}_{i \to v}) - \log p(\bm{x}) \approx \vg^\top(\ve(v) - \ve(x_i))$. The second term, the \textbf{distance term}, keeps the proposal local, penalizing distant candidates.\footnote{The two terms of \eqref{eq:bg-discrete-proposal} come from expanding the Gaussian log-density $-\tfrac{1}{2\alpha}\lVert \ve(v)-\ve(x_i)-\tfrac{\alpha}{2}\vg\rVert^2$, which yields the gradient-alignment term $\tfrac12\vg^\top(\ve(v)-\ve(x_i))$ and the distance term $-\tfrac{1}{2\alpha}\lVert\ve(v)-\ve(x_i)\rVert^2$, together with a term $-\tfrac{\alpha}{8}\lVert\vg\rVert^2$ that is constant in the candidate $v$ and is therefore discarded.} Everything the sampler knows about which token is good, as opposed to merely near, is contained in that gradient-alignment term, and if it is uninformative the proposal reduces to a distance-penalized random choice. The surrogate is therefore the sampler's entire claim to be gradient-guided rather than random, and Section~\ref{sec:results-linradius} tests it directly.

The \textbf{Continuous Langevin Sampler} (CLS), following the continuous-relaxation approach of methods such as MuCoLa \citep{kumar2022gradient}, takes the opposite route: it relaxes the state into the continuous embedding space $\R^D$, runs genuine Langevin dynamics there according to \eqref{eq:bg-langevin}, and projects back onto the nearest token embedding whenever a discrete sequence is required. The target density it navigates is therefore defined \emph{through} a \textbf{nearest-neighbour projection}, and that construction has a consequential geometry. The embedding space is partitioned into \textbf{Voronoi cells}, one per token; within a cell the projected token is constant, so the projected energy is constant and the cell is a flat plateau, and across a boundary the projected token changes abruptly and the energy jumps. The target is thus not the smooth surface Langevin theory assumes but a collection of plateaus separated by cliffs.

Writing $\mathrm{proj}_V(\bm{s})$ for the nearest-neighbour map that sends the continuous state at each masked position to the token whose embedding is closest in Euclidean distance, that target energy is
\begin{equation}
    \energy_{\text{CLS}}(\bm{s}) = -\log p_{\text{LM}}\bigl(\mathrm{proj}_V(\bm{s})\bigr),
    \label{eq:bg-cls-energy}
\end{equation}
the negative log-likelihood of the discrete sequence obtained by projecting the state back to tokens, to which the additive constraint $\lambda\, C$ of \eqref{eq:intro-energy} is appended when a constraint is imposed.

Two objects must be kept apart here, and the rest of the thesis depends on the distinction. The energy of \eqref{eq:bg-cls-energy} is piecewise constant and possesses no classical gradient, its derivative being zero inside every cell and undefined on the boundaries, so no sampler differentiates it. What the implemented sampler differentiates is a distinct \emph{differentiable pathway}: the log-likelihood evaluated with the continuous state supplied as the input embedding at the masked positions, while the projected token $\mathrm{proj}_V(\bm{s})$ supplies the prediction target.
% SOURCE: core/cls.py + core/base_sampler.get_gradient_and_log_joint + core/prep.py.
% Implementation detail: the log-likelihood is evaluated with the continuous state
% supplied as the input embedding at the masked positions (inputs_embeds) while the
% projected token proj_V(s) supplies the prediction target, so the within-cell surface
% is nearly, not exactly, flat; the load-bearing discontinuity is the target-token jump.
Inside a single cell the target token is fixed, so this pathway is a smooth function of $\bm{s}$ and its gradient is well defined; it is also, for the same reason, nearly but not exactly flat there, unlike the projected energy, which is exactly flat. Crossing a cell boundary changes the prediction target, so the pathway itself jumps, and the drift computed from it jumps with it. Every claim in this thesis about a discontinuous or non-Lipschitz drift attaches to this differentiable pathway \emph{across cell boundaries}, not to a classical gradient of the piecewise-constant energy of \eqref{eq:bg-cls-energy}, which does not exist. It is this arrangement, a pathway that is smooth within a cell and discontinuous across one, sampling a target that is flat within a cell and jumps across one, that puts the correction of Section~\ref{sec:bg-mh} under strain.

The arrangement has a direct consequence for the acceptance probability, derived once here. For a proposed move that stays within a single Voronoi cell the projected sequence, and therefore the energy of \eqref{eq:bg-cls-energy}, is unchanged, so the target ratio is at or near unity and the acceptance probability of \eqref{eq:bg-mh} is governed almost entirely by the proposal ratio. For the discrete sampler a within-cell move is the proposal of the same token, which is its own reverse, so the proposal ratio is one and the move is accepted with certainty. For the continuous sampler a within-cell move is still a continuous displacement whose reverse density is evaluated from the drift at the proposed point, which even inside the same cell need not make the return move likely, so within-cell acceptance is not protected. Section~\ref{sec:results-mh} reports the measured rates on both sides of this split.

\subsection{Score Matching and the Concrete Score}
\label{sec:bg-score}
\label{sec:bg-discrete-diffusion}

One further notion is needed before the experiments: the \textbf{score function}. For a differentiable density $p(\bm{s})$, the score is the gradient of its log-density, $\grad_{\bm{s}} \log p(\bm{s})$, and it is exactly the quantity the Langevin drift of \eqref{eq:bg-langevin} requires. \textbf{Score matching} \citep{vincent2011connection} fits a model directly to that quantity, and the observation behind modern score-based and diffusion generative models \citep{song2019generative, ho2020denoising} is that learning to denoise a corrupted input is equivalent to learning the score. A diffusion model, trained to reverse a gradual noising process, therefore acquires by construction a field of local directions that an autoregressive objective never asks for. That makes such a model one available source of a usable proposal, and Section~\ref{sec:results-diffusion} uses it as such; it is not the only possible source, and the results of Section~\ref{sec:results-mlm} show that it is not the necessary one either.

The construction does not transfer to text unchanged, because there is no gradient of a log-density with respect to a token index. The discrete analogue is the \textbf{concrete score} \citep{meng2022concrete}, which replaces the gradient by the collection of ratios $p(\bm{y})/p(\bm{x})$ between a sequence $\bm{x}$ and each neighbour $\bm{y}$ differing from it in one position. For every single-token move these ratios state how much more or less probable the data distribution makes it, which is the local directional signal a sampler needs in order to choose a substitution. Discrete diffusion models are trained to supply them. Their noising process corrupts a sequence through a Markov chain over the vocabulary; in the \textbf{absorbing} variant \citep{austin2021structured} each token is independently replaced, with a probability growing along the schedule, by a distinguished absorbing mask state, and the model is trained to reverse the process by predicting the clean token at each masked position conditioned on the surviving context on both sides. That reverse step is a bidirectional denoising query, so an absorbing diffusion model is trained to answer exactly the question an in-place revision poses: given a sequence with one position held out, which token belongs there. Two properties are therefore bundled together in such a model, the score objective and the bidirectional conditioning, and Section~\ref{sec:results-mlm} is the experiment that separates them.

Score-entropy discrete diffusion (SEDD) \citep{lou2024discrete}, the model used in Section~\ref{sec:results-diffusion}, learns the concrete score directly, fitting the ratios of the data distribution over the discrete state space with a loss adapted to non-negative ratio targets. It therefore provides, in a single forward pass, a trained estimate of the per-token probability ratios at exactly the positions a substitution could act on, and it shares the GPT-2 tokenizer, which is what lets it be dropped into this study's harness without retokenization.

\subsection{Amortized Inference with GFlowNets}
\label{sec:bg-gflownet}

Both samplers above require the gradient of the energy: the discrete sampler to score its proposals, the continuous sampler to compute its drift. An entirely different strategy is to give up on the gradient and instead \emph{learn} to sample. A \textbf{Generative Flow Network} \citep{bengio2021flow, bengio2023gflownet} treats generation as the sequential construction of an object and trains a policy so that the probability of terminating at a complete object is proportional to a non-negative reward $R(\bm{x})$, a unit of flow entering at the empty state and being split among the branches so that the flow arriving at each complete object matches its reward. Because flow is distributed in proportion to reward rather than concentrated on the single best object, the objects produced are \emph{diverse}, which is what distinguishes a GFlowNet from a reward-maximizing reinforcement-learning agent that would learn to produce only the one best object and would be useless as a sampler.

For text the object under construction is a sequence built from left to right, so the space of partial constructions is a tree with exactly one path from the empty state to any complete sequence. That structure simplifies the training objectives, which enforce a consistency condition on the flow along trajectories; those used in practice include trajectory balance \citep{malkin2022trajectory} and its sub-trajectory variants \citep{madan2023learning}, which improve credit assignment along long trajectories and are what the reference implementation this thesis builds on employs. One feature of the approach matters most here: the reward $R(\bm{x})$ is only ever \emph{evaluated}, never \emph{differentiated}. The policy learns from scalar scores of complete sequences, so the discontinuity and non-smoothness of the underlying landscape are invisible to it, and if the difficulty with energy-guided sampling were purely a property of the gradient, a method that never uses the gradient ought to escape it. Following \citet{hu2024amortizing}, the GFlowNet is realized as a parameter-efficient low-rank adapter \citep{hu2022lora} on the same frozen base model that supplies the energy for the Langevin samplers, so the two halves of the study operate on one shared model and their comparison is not confounded by any difference in the underlying network.
